\section{A line zero prevents closability of the specified smoothed Euler composition}
\label{sec:boundary-euler-nonclosability}

This is conditional on a zero on the chosen Mellin line. It
asserts neither the existence of an off-critical zero nor
failure of every possible transport. The sequential operator
domain is part of the statement.

Use (MA-MO.20), with \(0<\sigma<1\) and \(a=\sigma-1/2\).
Let \(A(t)=L(\sigma+it)E_S(\sigma+it)\) be the zeta or
primitive Dirichlet multiplier of MA-MO.21. More generally,
the argument applies to a function analytic near the real line,
not identically zero, whose derivatives of every fixed order
have polynomial growth there. Assume that \(A\) has a real
zero \(u_0\) of order \(m\ge1\). Let \(\phi\) be smooth,
one sufficiently far left and zero sufficiently far right,
as in the original source. Set
\[
 T_\phi=\mathcal F_+M_\phi\mathcal F_+^{-1},\qquad
 \mathscr C_A=C_c^\infty(\mathbb R\setminus\{A=0\}).
 \tag{MA-MO.69}
\]
The zero set is locally finite, so this core is dense in
ordinary \(L^2(\mathbb R,dt)\). For a nonzero
\(g\in C_c^\infty(\mathbb R)\), put
\[
 G(t)=\mathcal Mg(a+it),\qquad
 Q=M_AT_\phi M_{1/A},\qquad Q_g=M_GM_AT_\phi M_{1/A}
                  \quad\hbox{on }\mathscr C_A.
 \tag{MA-MO.70}
\]
Every displayed composition is sequentially defined on this
core: \(v/A\) is smooth and compactly supported,
\(T_\phi(v/A)\) is Schwartz, and multiplying by \(A\)
gives an \(L^2\) function. These are restrictions of the
literal maximal multiplier compositions.

\paragraph{MA-MO-006 (conditional nonclosability, including left smoothing).}
Under this line-zero hypothesis, neither \(Q\) nor \(Q_g\)
is closable on \(\mathscr C_A\). Neither is
\[
                 M_G(T_\phi-M_AT_\phi M_{1/A}).
 \tag{MA-MO.71}
\]
None therefore admits a closable extension agreeing on that
core. Under \(V_a\), (MA-MO.71) is the specified sequential
difference \(C_g(M_\phi-Z_{a,S}M_\phi J_{a,S})\).

\begin{proof}
Define
\[
 P(v)=\int e^{ivx}\phi'(x)\,dx,
             \qquad q(v)=\frac{iP(v)}{2\pi v}\quad(v\ne0).
 \tag{MA-MO.72}
\]
Then \(P\) is entire and Schwartz on the real line, and
\(P(0)=-1\). This is the off-diagonal kernel of MA-MO.22.
Distributional differentiation gives the full convolution
kernel \(\operatorname{pv}q+c\delta_0\) for some constant
\(c\): multiplying it by \(v\) gives \(iP(v)/(2\pi)\),
and the kernel of multiplication by \(v\) on distributions
is the span of \(\delta_0\). Its constant is immaterial.

For a smooth polynomially growing function \(B\), the
commutator on a compact smooth test has the ordinary kernel
\[
 \begin{aligned}
 [M_B,T_\phi]h(t)&=\int r_B(t,u)h(u)\,du,\\
 r_B(t,u)&=(B(t)-B(u))q(t-u)\quad(t\ne u),\\
 r_B(u,u)&=\frac{iP(0)}{2\pi}B'(u).
 \end{aligned}
 \tag{MA-MO.73}
\]
The delta term cancels and the difference removes the
principal-value singularity. For \(B=A\), both terms of
the commutator are in \(L^2\), by the domain observation
above.

The map \(u\mapsto r_A(\cdot,u)\) is continuous into
\(L^2(dt)\) for \(u\) near \(u_0\). To verify this,
restrict \(u\) to a fixed compact neighborhood. When
\(|t-u|\le1\), the difference quotient of \(A\) is
bounded by \(\sup|A'|\) on one fixed compact interval.
When \(|t-u|\ge1\), rapid decay of \(q(t-u)\) beats
the polynomial growth of \(A(t)-A(u)\), uniformly in
that neighborhood. These estimates give an integrable
square majorant. Pointwise continuity, including the
removable diagonal value, and dominated convergence give
the claimed \(L^2\) continuity.

Choose nonnegative \(h\in C_c^\infty((1,2))\) of integral
one and set
\[
 h_\epsilon(u)=\epsilon^{-1}h((u-u_0)/\epsilon),
                  \qquad v_\epsilon(u)=A(u)h_\epsilon(u).
 \tag{MA-MO.74}
\]
For sufficiently small positive \(\epsilon\), this support
avoids every real zero of \(A\), including \(u_0\).
Thus \(v_\epsilon\in\mathscr C_A\) and
\(M_{1/A}v_\epsilon=h_\epsilon\). Taylor expansion gives
\[
            \|v_\epsilon\|_2=O(\epsilon^{m-1/2})\to0.
 \tag{MA-MO.75}
\]
The exact commutator identity now yields
\[
 \begin{aligned}
 Qv_\epsilon
 &=T_\phi(Ah_\epsilon)+[M_A,T_\phi]h_\epsilon\\
 &\longrightarrow r_A(\cdot,u_0)=A(t)q(t-u_0)
                                      \quad\hbox{in }L^2.
 \end{aligned}
 \tag{MA-MO.76}
\]
The first term tends to zero by boundedness of \(T_\phi\).
The second is the approximate-identity average of the
proved \(L^2\)-continuous map. The last equality uses
\(A(u_0)=0\), with its removable value at \(u_0\).
Near that point, write \(A(t)=(t-u_0)^m a_0(t)\),
where \(a_0(u_0)\ne0\). Since \(P(0)=-1\), the limit
has a nonzero leading term of order \(m-1\). It is an
\(L^2\) function by the estimates above. The graph closure
therefore contains a nonzero vector above zero input,
which proves nonclosability of \(Q\).

The function \(G\) is bounded and Schwartz on the real
line and is the restriction of an entire function. It is
not identically zero, since Fourier uniqueness would
otherwise give \(e^{ax}g(x)=0\). Multiplying (MA-MO.76)
by this bounded function gives
\[
 Q_gv_\epsilon\longrightarrow G(t)A(t)q(t-u_0)\ne0
                                      \quad\hbox{in }L^2.
 \tag{MA-MO.77}
\]
If \(G(u_0)=0\), that zero has finite order, so the same
local expansion proves nonvanishing. Finally \(M_GT_\phi\)
is bounded and sends \(v_\epsilon\) to zero in norm.
Thus (MA-MO.71) has the negative nonzero limit and is
nonclosable too. Any extension containing the core has
the same graph sequence, proving the extension assertion.
\end{proof}

\paragraph{Domains and remaining scope.}
The inverse is applied only to tests supported away from
zeros, with every following domain verified. A factored
commutator involving \(J_{a,S}M_\phi\) would require an
additional inverse domain; no such equality of unbounded
compositions is used. The rational geometric comparison
MA-MO-004/005 has its own two-moment-zero smooth domain.
The present compact-frequency sequence is not asserted
to lie in that domain, so its multiplier nonclosability
is not imported into every original common-space model.
The theorem does not address the PR finite-jet quotient
returns or Gram-normalized maps with other domains.
The RS results on lines avoiding zeros retain their
hypotheses. On a line containing a zero, left multiplication
by nonzero compact smooth convolution does not make the
particular composition (MA-MO.70) closable.

\paragraph{Finite-length and nuclear-framework placement.}
The order \(m\) is the length of the local zero-divisor
module. Sequential domains and closability are additional
analytic structure, not consequences of that finite length.
The inherited cote~1 nuclear-trace locators, scans~079--080
and~111--112 through ALPOGE\_BRIDGES BR-TR-01/02, require
the relevant maps and their nuclear hypotheses before
traces or determinants are formed; a proposed closed
realization must also have a valid operator domain.
Neither finite-length \(K_0\) data nor arithmetic weights
repair the nonclosable composition proved here. No new
primary-source reading is claimed by this placement.
